\documentclass[11pt,a4paper]{article}
\usepackage[utf8]{inputenc}
\usepackage[T1]{fontenc}
\usepackage[brazil]{babel}
\usepackage[margin=2.5cm]{geometry}
\usepackage{titlesec}
\usepackage{enumitem}
\usepackage{booktabs}
\usepackage{tabularx}
\usepackage{array}
\usepackage{xcolor}
\usepackage{hyperref}
\usepackage{fancyhdr}

\definecolor{tfblue}{RGB}{20,40,70}
\definecolor{tfgold}{RGB}{200,150,40}

\titleformat{\section}{\normalfont\Large\bfseries\color{tfblue}}{\thesection}{0.6em}{}
\titleformat{\subsection}{\normalfont\large\bfseries\color{tfblue}}{\thesubsection}{0.6em}{}

\pagestyle{fancy}
\fancyhf{}
\fancyhead[L]{\small\textit{TrackFleet --- Qualidade}}
\fancyhead[R]{\small\thepage}
\renewcommand{\headrulewidth}{0.4pt}

\newcolumntype{L}{>{\raggedright\arraybackslash}X}
\newcommand{\dodbox}{$\square$}

\begin{document}

\begin{center}
{\LARGE\bfseries\color{tfblue} DOCUMENTO DE ESCOPO DO PROJETO}\\[2pt]
{\Large\bfseries TrackFleet --- Gestão de Rastreador de Automóveis}\\[4pt]
{\small\itshape Disciplina de Gerenciamento de Projetos $\cdot$ Framework PMBOK\textregistered\ 8\textsuperscript{a} Edição $\cdot$ Domínio: Escopo}\\[2pt]
{\small Integrantes: Enzo Allebrand e Kerlon Ribeiro (dois GPs) $\cdot$ 4\textsuperscript{o} semestre $\cdot$ Data: 20/08}
\end{center}

\vspace{2mm}
\noindent\textbf{\color{tfgold}\large Qualidade}

\vspace{1mm}
\noindent\textit{A terceira dimensão do escopo --- não um domínio à parte, mas um atributo integrante dele. A qualidade foca em atender às expectativas dos stakeholders e em cumprir os requisitos definidos tanto no Escopo do Produto quanto no Escopo do Projeto.}

\section{Requisitos de Qualidade do Produto}
\begin{itemize}[itemsep=2pt]
  \item Precisão e atualidade da localização exibida no mapa, com latência aceitável em relação à posição real do rastreador.
  \item Confiabilidade dos alertas de manutenção preventiva, com o mínimo possível de falsos positivos e falsos negativos.
  \item Usabilidade do painel: o gestor deve conseguir localizar um veículo e visualizar alertas pendentes em poucos cliques.
  \item Conformidade com a LGPD no tratamento de dados de localização e de motoristas (consentimento informado e minimização de dados coletados).
  \item Disponibilidade da plataforma durante toda a janela de testes com a transportadora piloto.
\end{itemize}

\section{Requisitos de Qualidade do Projeto (Processo)}
\begin{itemize}[itemsep=2pt]
  \item Todo código é revisado pelos dois GPs antes de ser integrado --- não há code review unilateral.
  \item Rotinas críticas (integração com a API do rastreador, cálculo de alertas) possuem testes automatizados.
  \item A documentação do projeto é atualizada a cada entrega, não deixada para o final.
\end{itemize}

\section{Definição de Pronto (DoD)}
A DoD é o conjunto de critérios que uma funcionalidade precisa cumprir para ser considerada apta ao uso pela transportadora piloto. Ela é definida \textbf{antes} de a funcionalidade começar a ser construída, para blindar o aceite e evitar a conversa em que a equipe diz que terminou e o patrocinador diz que não.

\begin{itemize}[itemsep=3pt]
  \item[\dodbox] Funcionalidade implementada conforme o requisito descrito no Escopo do Produto.
  \item[\dodbox] Código revisado por ambos os GPs (Enzo Allebrand e Kerlon Ribeiro).
  \item[\dodbox] Testada com dados reais ou simulados do rastreador GPS piloto.
  \item[\dodbox] Sem exposição indevida de dados pessoais (checagem de conformidade com a LGPD).
  \item[\dodbox] Validada e aceita pelos dois GPs; mudanças de escopo, prazo ou objetivo passam ainda pelo patrocinador, conforme a matriz RACI da Governança.
\end{itemize}

\section{Critérios de Aceite por Entregável}
\begin{center}
\begin{tabularx}{\textwidth}{L L}
\toprule
\textbf{Entregável} & \textbf{Critério de aceite} \\
\midrule
Monitoramento em tempo real (E1) & O gestor visualiza, no mapa, a posição atual de todos os veículos ativos com atraso perceptível mínimo \\
Gestão de manutenção preventiva (E2) & Um alerta é gerado automaticamente antes de o veículo atingir o limite de km/tempo configurado \\
Relatórios e indicadores (E4) & O painel exibe corretamente o número de veículos ativos e de alertas pendentes, com dados exportáveis \\
Piloto com a transportadora parceira & A transportadora utiliza a plataforma no período de teste e confirma, por escrito, que os indicadores refletem a realidade da frota \\
\bottomrule
\end{tabularx}
\end{center}

\section{Validar × Controlar o Escopo}
As duas atividades são complementares, mas distintas:
\begin{itemize}[itemsep=3pt]
  \item \textbf{Validar o Escopo} --- aceite formal, externo: o patrocinador (Diretoria da transportadora parceira) aprova cada marco relevante do projeto.
  \item \textbf{Controlar o Escopo} --- verificação interna, contínua: os dois GPs conferem, a cada check-in, se o que está sendo construído continua dentro da linha de base aprovada.
\end{itemize}
\noindent Um valida o \textbf{resultado}; o outro cuida do \textbf{processo}. Confundir os dois é um dos erros mais comuns de gestão de escopo.

\section{Prevenção de Scope Creep}
Os \textit{guardrails} já definidos no Documento de Governança sustentam também a qualidade do escopo:
\begin{itemize}[itemsep=2pt]
  \item Toda mudança relevante de escopo, prazo ou objetivo passa pelo patrocinador antes de ser aceita.
  \item Toda decisão importante é registrada no log de decisões do projeto, no momento em que é tomada.
  \item Nenhuma funcionalidade é considerada ``pronta'' sem passar pela DoD definida na Seção 3.
  \item Sem consenso entre os dois GPs, a decisão sobe direto ao patrocinador, que desempata.
\end{itemize}

\section{Métricas de Qualidade}
\begin{center}
\begin{tabularx}{\textwidth}{L l}
\toprule
\textbf{Métrica} & \textbf{Leitura} \\
\midrule
Taxa de defeitos encontrados no piloto & Menor é melhor \\
\% de funcionalidades aceitas sem retrabalho na primeira validação & Maior é melhor \\
\% de alertas de manutenção corretos (sem falso positivo/negativo) & Maior é melhor \\
\bottomrule
\end{tabularx}
\end{center}

\end{document}
